% Processing pipeline, one box per script (monochrome). Use with % Processing pipeline, one box per script (monochrome). Use with % Processing pipeline, one box per script (monochrome). Use with % Processing pipeline, one box per script (monochrome). Use with \input{figs/Ch4/processing-pipeline}.
% Needs tikz with arrows.meta and positioning (loaded in 0_main.tex). Font: TeX Gyre Heros (Helvetica clone).
\ifdefined\newfontfamily
    \ifdefined\pipelinefont\else
        \newfontfamily\pipelinefont{texgyreheros}[Extension=.otf, UprightFont=*-regular, BoldFont=*-bold]
    \fi
\else
    \providecommand{\pipelinefont}{\sffamily}
\fi
\begin{tikzpicture}[font=\pipelinefont\fontsize{8}{9.6}\selectfont,
  box/.style={draw=black, line width=0.4pt, align=center, minimum height=16mm, inner sep=0.6mm},
  ends/.style={draw=black, rounded corners=2mm, line width=0.4pt, align=center, minimum height=16mm, inner sep=1mm},
  stage/.style={font=\pipelinefont\fontsize{8}{9.6}\selectfont\bfseries, text height=1.6ex, text depth=0.4ex},
  flow/.style={-{Stealth[length=2mm]}, line width=0.4pt}]
  \node[ends, minimum width=14mm, text width=12.8mm] (n0) at (7.0mm,0) {Stereo\\recording};
  \node[box, minimum width=21mm, text width=19.8mm] (n1) at (27.5mm,0) {extract\_stereo\_\\frames.py\\(GUI)};
  \node[stage, above=1.2mm of n1] {Frame extraction};
  \node[box, minimum width=20mm, text width=18.8mm] (n2) at (51.0mm,0) {Calibration.py};
  \node[stage, above=1.2mm of n2] {Calibration};
  \node[box, minimum width=15mm, text width=13.8mm] (n3) at (71.5mm,0) {depth\_\\frames.py};
  \node[stage, above=1.2mm of n3] {Depth};
  \node[box, minimum width=23mm, text width=21.8mm] (n4) at (93.5mm,0) {tag\_poses.py\\colmap\_poses.py\\icp\_poses.py};
  \node[stage, above=1.2mm of n4] {Pose};
  \node[box, minimum width=21mm, text width=19.8mm] (n5) at (118.5mm,0) {tsdf\_fuse.py\\colmap\_mvs.py};
  \node[stage, above=1.2mm of n5] {Fusion};
  \node[box, minimum width=19mm, text width=17.8mm] (n6) at (141.5mm,0) {postprocess\_\\cloud.py};
  \node[stage, above=1.2mm of n6] {Clean-up};
  \node[ends, minimum width=11mm, text width=9.8mm] (n7) at (159.5mm,0) {Point\\cloud};
  \draw[flow] (n0) -- (n1);
  \draw[flow] (n1) -- (n2);
  \draw[flow] (n2) -- (n3);
  \draw[flow] (n3) -- (n4);
  \draw[flow] (n4) -- (n5);
  \draw[flow] (n5) -- (n6);
  \draw[flow] (n6) -- (n7);
\end{tikzpicture}
.
% Needs tikz with arrows.meta and positioning (loaded in 0_main.tex). Font: TeX Gyre Heros (Helvetica clone).
\ifdefined\newfontfamily
    \ifdefined\pipelinefont\else
        \newfontfamily\pipelinefont{texgyreheros}[Extension=.otf, UprightFont=*-regular, BoldFont=*-bold]
    \fi
\else
    \providecommand{\pipelinefont}{\sffamily}
\fi
\begin{tikzpicture}[font=\pipelinefont\fontsize{8}{9.6}\selectfont,
  box/.style={draw=black, line width=0.4pt, align=center, minimum height=16mm, inner sep=0.6mm},
  ends/.style={draw=black, rounded corners=2mm, line width=0.4pt, align=center, minimum height=16mm, inner sep=1mm},
  stage/.style={font=\pipelinefont\fontsize{8}{9.6}\selectfont\bfseries, text height=1.6ex, text depth=0.4ex},
  flow/.style={-{Stealth[length=2mm]}, line width=0.4pt}]
  \node[ends, minimum width=14mm, text width=12.8mm] (n0) at (7.0mm,0) {Stereo\\recording};
  \node[box, minimum width=21mm, text width=19.8mm] (n1) at (27.5mm,0) {extract\_stereo\_\\frames.py\\(GUI)};
  \node[stage, above=1.2mm of n1] {Frame extraction};
  \node[box, minimum width=20mm, text width=18.8mm] (n2) at (51.0mm,0) {Calibration.py};
  \node[stage, above=1.2mm of n2] {Calibration};
  \node[box, minimum width=15mm, text width=13.8mm] (n3) at (71.5mm,0) {depth\_\\frames.py};
  \node[stage, above=1.2mm of n3] {Depth};
  \node[box, minimum width=23mm, text width=21.8mm] (n4) at (93.5mm,0) {tag\_poses.py\\colmap\_poses.py\\icp\_poses.py};
  \node[stage, above=1.2mm of n4] {Pose};
  \node[box, minimum width=21mm, text width=19.8mm] (n5) at (118.5mm,0) {tsdf\_fuse.py\\colmap\_mvs.py};
  \node[stage, above=1.2mm of n5] {Fusion};
  \node[box, minimum width=19mm, text width=17.8mm] (n6) at (141.5mm,0) {postprocess\_\\cloud.py};
  \node[stage, above=1.2mm of n6] {Clean-up};
  \node[ends, minimum width=11mm, text width=9.8mm] (n7) at (159.5mm,0) {Point\\cloud};
  \draw[flow] (n0) -- (n1);
  \draw[flow] (n1) -- (n2);
  \draw[flow] (n2) -- (n3);
  \draw[flow] (n3) -- (n4);
  \draw[flow] (n4) -- (n5);
  \draw[flow] (n5) -- (n6);
  \draw[flow] (n6) -- (n7);
\end{tikzpicture}
.
% Needs tikz with arrows.meta and positioning (loaded in 0_main.tex). Font: TeX Gyre Heros (Helvetica clone).
\ifdefined\newfontfamily
    \ifdefined\pipelinefont\else
        \newfontfamily\pipelinefont{texgyreheros}[Extension=.otf, UprightFont=*-regular, BoldFont=*-bold]
    \fi
\else
    \providecommand{\pipelinefont}{\sffamily}
\fi
\begin{tikzpicture}[font=\pipelinefont\fontsize{8}{9.6}\selectfont,
  box/.style={draw=black, line width=0.4pt, align=center, minimum height=16mm, inner sep=0.6mm},
  ends/.style={draw=black, rounded corners=2mm, line width=0.4pt, align=center, minimum height=16mm, inner sep=1mm},
  stage/.style={font=\pipelinefont\fontsize{8}{9.6}\selectfont\bfseries, text height=1.6ex, text depth=0.4ex},
  flow/.style={-{Stealth[length=2mm]}, line width=0.4pt}]
  \node[ends, minimum width=14mm, text width=12.8mm] (n0) at (7.0mm,0) {Stereo\\recording};
  \node[box, minimum width=21mm, text width=19.8mm] (n1) at (27.5mm,0) {extract\_stereo\_\\frames.py\\(GUI)};
  \node[stage, above=1.2mm of n1] {Frame extraction};
  \node[box, minimum width=20mm, text width=18.8mm] (n2) at (51.0mm,0) {Calibration.py};
  \node[stage, above=1.2mm of n2] {Calibration};
  \node[box, minimum width=15mm, text width=13.8mm] (n3) at (71.5mm,0) {depth\_\\frames.py};
  \node[stage, above=1.2mm of n3] {Depth};
  \node[box, minimum width=23mm, text width=21.8mm] (n4) at (93.5mm,0) {tag\_poses.py\\colmap\_poses.py\\icp\_poses.py};
  \node[stage, above=1.2mm of n4] {Pose};
  \node[box, minimum width=21mm, text width=19.8mm] (n5) at (118.5mm,0) {tsdf\_fuse.py\\colmap\_mvs.py};
  \node[stage, above=1.2mm of n5] {Fusion};
  \node[box, minimum width=19mm, text width=17.8mm] (n6) at (141.5mm,0) {postprocess\_\\cloud.py};
  \node[stage, above=1.2mm of n6] {Clean-up};
  \node[ends, minimum width=11mm, text width=9.8mm] (n7) at (159.5mm,0) {Point\\cloud};
  \draw[flow] (n0) -- (n1);
  \draw[flow] (n1) -- (n2);
  \draw[flow] (n2) -- (n3);
  \draw[flow] (n3) -- (n4);
  \draw[flow] (n4) -- (n5);
  \draw[flow] (n5) -- (n6);
  \draw[flow] (n6) -- (n7);
\end{tikzpicture}
.
% Needs tikz with arrows.meta and positioning (loaded in 0_main.tex). Font: TeX Gyre Heros (Helvetica clone).
\ifdefined\newfontfamily
    \ifdefined\pipelinefont\else
        \newfontfamily\pipelinefont{texgyreheros}[Extension=.otf, UprightFont=*-regular, BoldFont=*-bold]
    \fi
\else
    \providecommand{\pipelinefont}{\sffamily}
\fi
\begin{tikzpicture}[font=\pipelinefont\fontsize{8}{9.6}\selectfont,
  box/.style={draw=black, line width=0.4pt, align=center, minimum height=16mm, inner sep=0.6mm},
  ends/.style={draw=black, rounded corners=2mm, line width=0.4pt, align=center, minimum height=16mm, inner sep=1mm},
  stage/.style={font=\pipelinefont\fontsize{8}{9.6}\selectfont\bfseries, text height=1.6ex, text depth=0.4ex},
  flow/.style={-{Stealth[length=2mm]}, line width=0.4pt}]
  \node[ends, minimum width=14mm, text width=12.8mm] (n0) at (7.0mm,0) {Stereo\\recording};
  \node[box, minimum width=21mm, text width=19.8mm] (n1) at (27.5mm,0) {extract\_stereo\_\\frames.py\\(GUI)};
  \node[stage, above=1.2mm of n1] {Frame extraction};
  \node[box, minimum width=20mm, text width=18.8mm] (n2) at (51.0mm,0) {Calibration.py};
  \node[stage, above=1.2mm of n2] {Calibration};
  \node[box, minimum width=15mm, text width=13.8mm] (n3) at (71.5mm,0) {depth\_\\frames.py};
  \node[stage, above=1.2mm of n3] {Depth};
  \node[box, minimum width=23mm, text width=21.8mm] (n4) at (93.5mm,0) {tag\_poses.py\\colmap\_poses.py\\icp\_poses.py};
  \node[stage, above=1.2mm of n4] {Pose};
  \node[box, minimum width=21mm, text width=19.8mm] (n5) at (118.5mm,0) {tsdf\_fuse.py\\colmap\_mvs.py};
  \node[stage, above=1.2mm of n5] {Fusion};
  \node[box, minimum width=19mm, text width=17.8mm] (n6) at (141.5mm,0) {postprocess\_\\cloud.py};
  \node[stage, above=1.2mm of n6] {Clean-up};
  \node[ends, minimum width=11mm, text width=9.8mm] (n7) at (159.5mm,0) {Point\\cloud};
  \draw[flow] (n0) -- (n1);
  \draw[flow] (n1) -- (n2);
  \draw[flow] (n2) -- (n3);
  \draw[flow] (n3) -- (n4);
  \draw[flow] (n4) -- (n5);
  \draw[flow] (n5) -- (n6);
  \draw[flow] (n6) -- (n7);
\end{tikzpicture}
